\documentclass[../../main.tex]{subfiles}

\begin{document}

\chapter{Kriptografi Kunci-Publik: Algoritma ElGamal}
\label{chap:elgamal}

\begin{learningoutcome}
Setelah mempelajari bab ini, mahasiswa diharapkan mampu:
\begin{itemize}
    \item Menjelaskan konsep dasar algoritma ElGamal dan hubungannya dengan masalah logaritma diskrit (Sub-CPMK 1.3).
    \item Menguraikan prosedur pembangkitan pasangan kunci publik dan privat pada sistem ElGamal (Sub-CPMK 1.3).
    \item Menerapkan rumus enkripsi dan dekripsi ElGamal untuk mengamankan pesan (Sub-CPMK 1.3).
    \item Menganalisis properti keamanan ElGamal, termasuk sifat \textit{probabilistic encryption} (Sub-CPMK 1.3).
    \item Membandingkan mekanisme kerja ElGamal dengan RSA dan Diffie-Hellman (Sub-CPMK 1.3).
\end{itemize}
\end{learningoutcome}

\subfile{section-01}
\subfile{section-02}
\subfile{section-03}
\subfile{section-04}
\section{Contoh Terhitung}
\label{sec:contoh-elgamal}

Seluruh contoh pada bagian ini memakai aritmetika modular yang telah dibahas
pada Bab~\ref{chap:landasan-matematika}. Contoh pertama sengaja dipilih sekecil
mungkin agar setiap langkah dapat diperiksa dengan tangan; contoh berikutnya
memakai modulus yang lebih besar dan pesan yang lebih panjang untuk
memperlihatkan bahwa prosedurnya sama sekali tidak berubah ketika ukurannya
bertambah. Dua contoh terakhir memperlihatkan dua sifat yang hanya muncul pada
pemakaian yang lebih rumit: sifat homomorfik dan akibat penggunaan ulang kunci
efemeral.

\begin{example}[Membandingkan akar primitif dan bukan akar primitif]
\label{ex:elgamal-akar}
Misalkan $p = 7$. Bilangan $a = 3$ adalah akar primitif dari $7$, sebab
perpangkatannya menghasilkan keenam residu bukan-nol tepat sekali
masing-masing:
\[ 3^{1}, 3^{2}, 3^{3}, 3^{4}, 3^{5}, 3^{6} \equiv 3, 2, 6, 4, 5, 1 \pmod{7} . \]
Sebaliknya $a = 2$ bukan akar primitif, sebab
\[ 2^{1}, 2^{2}, 2^{3}, 2^{4}, 2^{5}, 2^{6} \equiv 2, 4, 1, 2, 4, 1 \pmod{7} , \]
sehingga orderya hanya $3$ dan tiga residu tidak pernah muncul sebagai kunci
publik. Kedua baris itu diperlihatkan berdampingan pada
Tabel~\ref{tab:akar-primitif-7}.

Untuk modulus yang sedikit lebih besar, logaritma diskret masih dapat
diselesaikan dengan mencoba satu per satu. Karena $7$ adalah akar primitif dari
$p = 41$, persamaan
\[ 7^{x} \equiv 15 \pmod{41} \]
dapat diselesaikan dengan menghitung $7^{1} = 7$, $7^{2} = 49 \equiv 8$,
$7^{3} = 343 \equiv 15$; jadi $x = 3$. Pada $p = 7$ jumlah akar primitifnya
adalah $\varphi(6) = 2$, yaitu $3$ dan $5$; pada $p = 41$ jumlahnya
$\varphi(40) = 16$.
\end{example}

\begin{example}[Pembangkitan pasangan kunci ElGamal]
\label{ex:elgamal-kunci}
Alice memilih bilangan prima $p = 2273$ dan akar primitif $g = 3$ (lihat
Gambar~\ref{fig:simulasi-enkripsi-elgamal}). Ia menetapkan kunci privat
$x = 243$, yang memenuhi syarat $2 \le 243 \le 2271$, lalu menghitung kunci
publiknya dengan Persamaan~\eqref{eq:elgamal-kunci}:
\[ y = g^{x} \bmod p = 3^{243} \bmod 2273 = 461 . \]
Kunci publik Alice adalah $(y, g, p) = (461, 3, 2273)$, sedangkan kunci
privatnya $(x, p) = (243, 2273)$.

Perhatikan bahwa kunci privat Alice hanya sebuah eksponen, bukan sepasang
bilangan prima seperti pada RSA. Yang perlu ia jaga hanyalah nilai $243$ itu,
sebab dari situlah seluruh keamanan pesannya bersumber. Kunci publiknya boleh
disiarkan melalui papan pengumuman tanpa pengamanan tambahan; nilai $461$,
$3$, dan $2273$ tidak mengungkapkan apa pun tentang $243$ selain bahwa
menemukannya menuntut penyelesaian persoalan logaritma diskret.
\end{example}

\begin{example}[Enkripsi dan dekripsi satu blok]
\label{ex:elgamal-enkripsi}
Bob mengirim pesan $m = 700$ kepada Alice, dan nilai itu memang berada di dalam
selang $[0, 2272]$. Ia memilih kunci efemeral $k = 1463$, yang juga berada di
dalam selang $[1, 2271]$.

\begin{enumerate}
    \item Enkripsi oleh Bob memakai kunci publik Alice
          $(y, g, p) = (461, 3, 2273)$:
    \begin{align*}
    a &= g^{k} \bmod p = 3^{1463} \bmod 2273 = 1439, \\
    b &= y^{k} \cdot m \bmod p = 461^{1463} \times 700 \bmod 2273 = 74 .
    \end{align*}
    Jadi cipherteks yang dikirim adalah $(a, b) = (1439, 74)$.
    \item Dekripsi oleh Alice dengan kunci privatnya $x = 243$ memakai
          Persamaan~\eqref{eq:elgamal-invers} dan
          Persamaan~\eqref{eq:elgamal-dekripsi}:
    \[
    \begin{aligned}
    \left(a^{x}\right)^{-1} &\equiv a^{\,p-1-x} \pmod{p} \\
                            &= 1439^{\,2273-1-243} \bmod 2273
                             = 1439^{\,2029} \bmod 2273 = 1791,
    \end{aligned}
    \]
    \[ m = b \cdot \left(a^{x}\right)^{-1} \bmod p
          = 74 \times 1791 \bmod 2273 = 700 . \]
\end{enumerate}

Plainteks $m = 700$ berhasil dipulihkan dengan tepat. Perhatikan bahwa eksponen
$2029$ pada langkah dekripsi berasal dari $p - 1 - x = 2273 - 1 - 243$, yaitu
jalan pintas untuk menghitung invers tanpa algoritma Euclidean.

Karena $k$ dipilih secara acak, enkripsi ulang pesan $m = 700$ yang sama dengan
$k$ yang berbeda akan menghasilkan pasangan $(a, b)$ yang berbeda pula. Inilah
yang dimaksud dengan enkripsi probabilistik pada
Subbagian~\ref{subsec:k-ulang} di muka.
\end{example}

\begin{example}[Pesan multi-blok ``HALO'' dengan dua kunci efemeral berbeda]
\label{ex:elgamal-halo}
Alice tetap memakai pasangan kunci pada Contoh~\ref{ex:elgamal-kunci}, yaitu
kunci publik $(y, g, p) = (461, 3, 2273)$. Kali ini Bob ingin mengirim kata
``HALO'' kepada Alice. Dengan pengodean $A = 00, B = 01, \dots, Z = 25$, setiap
huruf berubah menjadi dua digit, sebagaimana diperlihatkan pada
Tabel~\ref{tab:pengodean-halo}.

\begin{table}[htbp]
    \centering
    \begin{tabular}{@{}l c c c c@{}}
    \toprule
    Huruf & H & A & L & O \\
    \midrule
    Kode  & 07 & 00 & 11 & 14 \\
    \bottomrule
    \end{tabular}
    \caption{Pengodean kata ``HALO'' menjadi bilangan dengan aturan $A = 00$
    sampai $Z = 25$, sehingga pesannya menjadi $m = 07001114$}
    \label{tab:pengodean-halo}
\end{table}

Rangkaian angka itu menjadi $m = 07001114$. Karena $m$ jauh lebih besar
daripada modulus $p - 1 = 2272$, Bob harus memecahnya menjadi blok-blok yang
masing-masing lebih kecil daripada modulus. Dengan panjang blok empat digit
diperoleh $m_{1} = 0700$ dan $m_{2} = 1114$; keduanya memenuhi syarat
$0 \le m_{i} \le 2272$.

\begin{enumerate}
    \item \textbf{Enkripsi $m_{1} = 0700$ dengan $k = 1463$.} Nilai ini
          kebetulan sama dengan kunci efemeral pada
          Contoh~\ref{ex:elgamal-enkripsi}, sehingga hasilnya juga sama:
    \begin{align*}
    a_{1} &= 3^{1463} \bmod 2273 = 1439, \\
    b_{1} &= 461^{1463} \times 700 \bmod 2273 = 74 .
    \end{align*}
    Cipherteks untuk blok pertama adalah $c_{1} = (1439, 74)$.
    \item \textbf{Enkripsi $m_{2} = 1114$ dengan $k = 2001$.} Kali ini Bob
          memilih kunci efemeral yang \emph{berbeda}, sebagaimana mestinya:
    \begin{align*}
    a_{2} &= 3^{2001} \bmod 2273 = 1220, \\
    b_{2} &= 461^{2001} \times 1114 \bmod 2273 = 1682 .
    \end{align*}
    Cipherteks untuk blok kedua adalah $c_{2} = (1220, 1682)$.
\end{enumerate}

Bob mengirimkan kedua pasangan itu kepada Alice. Perhatikan bahwa ukuran
cipherteksnya kini dua kali lipat: pesan semula delapan digit
($07001114$), sedangkan cipherteksnya empat bilangan sehingga menjadi enam
belas digit.

Alice mendekripsinya dengan kunci privat $x = 243$.

\begin{enumerate}
    \item Untuk $c_{1} = (1439, 74)$:
    \[ \left(a_{1}^{x}\right)^{-1} = 1439^{\,2029} \bmod 2273 = 1791 , \]
    \[ m_{1} = 74 \times 1791 \bmod 2273 = 700 = 0700 . \]
    \item Untuk $c_{2} = (1220, 1682)$:
    \[ \left(a_{2}^{x}\right)^{-1} = 1220^{\,2029} \bmod 2273 = 1125 , \]
    \[ m_{2} = 1682 \times 1125 \bmod 2273 = 1114 . \]
\end{enumerate}

Hasilnya adalah $m_{1} m_{2} = 07001114$, yang dengan pengodean pada
Tabel~\ref{tab:pengodean-halo} berbunyi ``HALO'' --- sama persis dengan pesan
yang dikirim Bob.

Contoh ini memperlihatkan dua hal sekaligus. Pertama, pemecahan menjadi blok
memang diperlukan begitu pesan melebihi modulus, persis seperti pada RSA.
Kedua, kedua blok memakai kunci efemeral yang berbeda. Andaikan Bob memakai
$k = 1463$ untuk keduanya, komponen pertama dari kedua cipherteks akan sama
persis, dan itulah tanda yang membuka serangan pada
Contoh~\ref{ex:elgamal-k-ulang}.
\end{example}

\begin{example}[Parameter kedua: modulus 2357 dengan basis 2]
\label{ex:elgamal-2357}
Agar tidak terkesan bahwa hasilnya bergantung pada satu himpunan parameter
tertentu, contoh berikut memakai bilangan prima yang lain, basis yang lain,
serta kunci privat yang panjang. Misalkan $p = 2357$ dan $g = 2$, yang
keduanya dapat diperiksa: $2357$ prima dan orde $2$ modulo $2357$ sama dengan
$2356$, sehingga $2$ benar-benar akar primitif.

Bob memilih kunci privat $x = 1751$ dan menghitung
\[ y = 2^{1751} \bmod 2357 = 1185 . \]
Kunci publiknya adalah $(y, g, p) = (1185, 2, 2357)$.

Alice ingin mengirim $m = 2035$ dan memilih $k = 1520$. Enkripsinya:
\begin{align*}
a &= 2^{1520} \bmod 2357 = 1430, \\
b &= 1185^{1520} \times 2035 \bmod 2357 = 697 .
\end{align*}
Cipherteksnya $(a, b) = (1430, 697)$. Bob mendekripsinya dengan $x = 1751$:
\[ \left(a^{x}\right)^{-1} = 1430^{\,2357-1-1751} \bmod 2357 = 1430^{\,605} \bmod 2357 = 872 , \]
\[ m = 697 \times 872 \bmod 2357 = 2035 . \]
Plainteks $2035$ pulih kembali.

Perhatikan eksponen $605$ pada langkah dekripsi. Bila ditulis
$p - 1 - x = 2357 - 1 - 1751$, hasilnya memang $605$. Sumber asli contoh ini
menuliskannya sebagai $2^{1751} \bmod 2353$ dan bukan $\bmod 2357$; modulus
$2353$ jelas keliru, sebab tidak sama dengan $p$ yang telah ditetapkan, dan
memang $2^{1751} \bmod 2353 = 1008 \ne 1185$. Nilai yang benar diperoleh dengan
modulus $2357$.
\end{example}

\begin{example}[Membongkar pesan dengan dua kunci efemeral yang sama]
\label{ex:elgamal-k-ulang}
Sekarang diperlihatkan apa yang terjadi bila aturan pada
Contoh~\ref{ex:elgamal-halo} dilanggar. Andaikan Bob memakai kunci efemeral
$k = 1463$ untuk kedua bloknya, yaitu $m_{1} = 700$ dan $m_{2} = 1114$, pada
kunci publik Alice $(y, g, p) = (461, 3, 2273)$.

Karena $a = g^{k}$ tidak bergantung pada pesannya, komponen pertama kedua
cipherteks sama persis:
\[ a_{1} = a_{2} = 3^{1463} \bmod 2273 = 1439 . \]
Komponen keduanya berbeda, sebab pesannya berbeda:
\begin{align*}
b_{1} &= 461^{1463} \times 700 \bmod 2273 = 74, \\
b_{2} &= 461^{1463} \times 1114 \bmod 2273 = 988 .
\end{align*}

Sebuah penyadap yang memergoki kedua pasangan itu tidak perlu menghitung
logaritma diskret apa pun. Ia cukup mengikuti
Persamaan~\eqref{eq:elgamal-k-ulang}:
\[ \frac{b_{1}}{b_{2}} \equiv b_{1} \cdot b_{2}^{-1} \equiv 74 \cdot 988^{-1}
   \equiv 474 \pmod{2273} , \]
dan memang
$\frac{m_{1}}{m_{2}} \equiv 700 \cdot 1114^{-1} \equiv 474 \pmod{2273}$,
sehingga kedua rasio itu cocok.

Bila penyadap mengetahui bahwa blok pertama berbunyi $0700$ --- misalnya
karena setiap pesan pada sistem itu diawali dengan penanda yang tetap --- maka
ia langsung memperoleh blok kedua dengan
Persamaan~\eqref{eq:elgamal-k-ulang-m2}:
\[ m_{2} = m_{1} \cdot b_{2} \cdot b_{1}^{-1}
          = 700 \times 988 \times 74^{-1} \equiv 1114 \pmod{2273} . \]
Blok kedua terbaca sebagai ``$1114$'', yaitu ``LO''. Bila penanda awalnya tidak
diketahui, penyadap masih dapat mencoba setiap kandidat $m_{1}$ yang masuk
akal, menghitung $m_{2}$ yang berpasangan, dan memilih pasangan yang keduanya
bermakna --- pekerjaan yang jauh lebih ringan daripada menyelesaikan
logaritma diskret.

Sekali salah satu pasangan plainteks diperoleh, penyadap dapat menghitung
$y^{k} = b_{1} \cdot m_{1}^{-1} \bmod p = 74 \cdot 700^{-1} \bmod 2273$ dan
memakainya untuk membuka setiap cipherteks lain yang memakai $k$ yang sama.
Kunci privat Alice $x = 243$ sendiri tetap tidak terungkap; yang bocor adalah
isi pesannya. Bagi Alice, akibatnya tidak berbeda.
\end{example}


\begin{obeactivity}
\textbf{Aktivitas 12.1: Simulasi ElGamal dengan Bilangan Kecil}
1. Gunakan parameter publik $p = 2357, g = 2$.
2. Bob memilih kunci privat $x = 1751$, hitung kunci publik $y$.
3. Alice ingin mengirim pesan $m = 2035$ dengan memilih $k = 1520$. Hitunglah pasangan cipherteks $(a, b)$.
4. Bob melakukan dekripsi menggunakan kunci privatnya. Verifikasi apakah hasilnya kembali menjadi $m = 2035$.
\end{obeactivity}

Langkah-langkah pada Aktivitas 12.1 dapat dijalankan langsung melalui program
pada Listing~\ref{lst:elgamal}. Program itu membangkitkan kunci publik dari
kunci privat, mengenkripsi pesan menjadi pasangan $(a, b)$, lalu memulihkan
pesan semula.

\lstinputlisting[
    language=Python,
    caption={ElGamal: pembangkitan kunci, enkripsi, dan dekripsi},
    label={lst:elgamal}
]{\codefile{python/elgamal.py}}

Program yang sama tersedia dalam C++ (\texttt{code/cpp/elgamal.cpp}) dan MATLAB
(\texttt{code/matlab/elgamal.m}).


\begin{obereflection}
\textbf{Latihan 12.1}
1. Apa yang terjadi jika pengirim menggunakan nilai $k$ yang sama untuk dua pesan yang berbeda? Analisis risikonya terhadap keamanan pesan.
2. Jelaskan mengapa pemilihan $g$ sebagai akar primitif dari $p$ sangat penting dalam algoritma ElGamal.
3. Mengapa ukuran cipherteks ElGamal selalu dua kali lipat dari ukuran plainteksnya?
4. Diberikan parameter $p = 2357$, $g = 2$, dan kunci privat $x = 900$. Hitunglah
   kunci publik $y$, lalu enkripsikan $m = 1234$ dengan $k = 1000$. Periksa bahwa
   dekripsinya memang mengembalikan $1234$.
5. Cipherteks $(a_{1}, b_{1}) = (1439, 74)$ adalah hasil enkripsi $m_{1} = 700$
   pada modulus $p = 2273$. Bila penyerang memperoleh cipherteks lain
   $(a_{2}, b_{2}) = (1220, 1682)$ untuk pesan yang tidak diketahui, dapatkah ia
   memperoleh pesan kedua itu? Jelaskan langkah yang harus dilakukannya, atau
   jelaskan mengapa ia tidak dapat melakukannya.
6. Mengapa invers $\left(a^{x}\right)^{-1}$ dapat dihitung sebagai
   $a^{\,p-1-x} \bmod p$? Sebutkan teorema yang mendasarinya, dan syarat apa yang
   harus dipenuhi oleh $a$ agar teorema itu berlaku.
7. Jelaskan mengapa kunci privat $x = 1$ dan $x = p-1$ dikecualikan dari selang
   pemilihan kunci. Tunjukkan apa yang terjadi pada kunci publik $y$ dalam kedua
   kasus itu.
8. Misalkan $p = 2273$ dan pesan yang dikirim adalah rangkaian delapan digit
   $07001114$. Bandingkan pemecahan menjadi blok empat digit dengan pemecahan
   menjadi blok tiga digit. Berapa blok yang dihasilkan masing-masing cara, dan
   masalah apa yang muncul pada pemecahan tiga digit?
\end{obereflection}

\begin{obereflection}
\textbf{Latihan 12.2}
1. Sistem ElGamal memakai $k$ baru untuk setiap blok pesan, sehingga satu blok
   memerlukan dua perpangkatkan modular. Bandingkan biaya itu dengan RSA, yang
   hanya memerlukan satu perpangkatkan per blok, dan jelaskan mengapa ElGamal
   tetap dipandang layak dipakai.
2. Sebuah berkas multimedia berukuran 1 MB hendak dikirim dengan aman. Mengapa
   tidak tepat mengenkripsinya langsung dengan ElGamal? Rancanglah skema
   kriptografi hibrida yang memakai ElGamal untuk membungkus kunci sesi dan AES
   untuk isinya, lalu sebutkan besaran apa saja yang melewati saluran.
3. Tunjukkan bahwa ElGamal bersifat homomorfik multiplikatif. Perlihatkan
   langkahnya untuk $m_{1} = 700$ dan $m_{2} = 1114$ pada modulus $2273$, lalu
   jelaskan mengapa sifat itu berbahaya bagi keutuhan pesan.
4. OpenPGP pernah mencantumkan ElGamal sebagai algoritma kunci publik nomor 16
   dengan keterangan \textit{encrypt-only}. Jelaskan mengapa keterangan itu
   diberikan, dan mengapa RFC~9580 (2024) kemudian melarang pembangkitan kunci
   ElGamal yang baru.
5. Bandingkan ElGamal dengan protokol Diffie--Hellman dari sudut kebutuhan
   interaksi. Dalam keadaan apa keunggulan ElGamal itu menjadi menentukan, dan
   dalam keadaan apa keunggulan itu tidak berguna?
6. Jelaskan mengapa algoritma Shor mengancam baik RSA maupun ElGamal, sekalipun
   keduanya bersandar pada persoalan yang berbeda. Apa yang harus dilakukan
   sebuah lembaga yang hari ini menyimpan data rahasia jangka panjang?
\end{obereflection}

\section{Rangkuman}
\begin{itemize}
    \item ElGamal (Taher ElGamal, 1985) adalah sistem kunci-publik yang keamanannya bersandar pada persoalan logaritma diskrit, erat terkait dengan Diffie-Hellman.
    \item Pembangkitan kunci: pilih prima $p$ dan akar primitif $g$, pilih kunci privat $x$ dengan $2 \le x \le p-2$, lalu hitung kunci publik $y = g^x \bmod p$.
    \item Enkripsi bersifat probabilistik: pilih bilangan acak $k$, hitung $a = g^k \bmod p$ dan $b = y^k \cdot m \bmod p$, sehingga cipherteksnya berupa pasangan $(a, b)$.
    \item Dekripsi: $m = b \cdot (a^x)^{-1} \bmod p$, dengan $(a^x)^{-1} \equiv a^{p-1-x} \pmod p$.
    \item Contoh: $p = 2273$, $g = 3$, dan $x = 243$ memberi $y = 461$; pesan $m = 700$ dengan $k = 1463$ menjadi $(a, b) = (1439, 74)$ dan kembali menjadi $700$ setelah dekripsi.
    \item Karena $k$ berubah setiap sesi, plainteks yang sama menghasilkan cipherteks berbeda; hal ini mencegah serangan \textit{codebook} dan analisis statistik.
    \item Kelemahannya adalah ekspansi cipherteks (dua kali panjang plainteks); dibanding RSA, ElGamal bersifat probabilistik dan mendasari skema tanda tangan DSA.
\end{itemize}


\begin{obeassessment}
\textbf{Kuis Bab 12}
1. Jelaskan perbedaan mendasar antara masalah logaritma diskrit dan masalah pemfaktoran bilangan bulat.
2. Uraikan langkah-langkah pembangkitan kunci pada sistem ElGamal.
3. Mengapa ElGamal disebut sebagai sistem enkripsi probabilistik?
\end{obeassessment}
\begin{competencychecklist}
\begin{tabularx}{\linewidth}{|X|c|}
\hline
Indikator Kompetensi & Tercapai \\
\hline
Saya dapat menjelaskan konsep logaritma diskrit dan akar primitif & [ ] \\
\hline
Saya dapat menghitung kunci publik dan privat ElGamal & [ ] \\
\hline
Saya dapat melakukan proses enkripsi dan dekripsi ElGamal & [ ] \\
\hline
Saya dapat menganalisis efek penggunaan kunci efemeral $k$ & [ ] \\
\hline
Saya dapat membandingkan ElGamal dengan RSA & [ ] \\
\hline
\end{tabularx}
\end{competencychecklist}


\end{document}
